(1) 取 \(y(x,t)=X(x)e^{-i\omega t}\)，则
\[
-\mu\omega^2X=\frac d{dx}\left[T_0(1+\alpha x)^2X'\right],\qquad X(0)=X(L)=0.
\]
因为张力位置依赖，不能直接把波速当常数而写 \(X''+k^2X=0\)。

(2) 用 \(z=1+\alpha x\)，空间方程变为
\[
z^2X_{zz}+2zX_z+\frac{\mu\omega^2}{T_0\alpha^2}X=0.
\]
这是 Euler–Cauchy 方程；试 \(X=z^\nu\) 得 \(\nu(\nu+1)+\mu\omega^2/(T_0\alpha^2)=0\)。对可满足两个固定端的振动模，写 \(\nu=-1/2\pm iq\)，因此
\[
X(z)=z^{-1/2}[A\cos(q\ln z)+B\sin(q\ln z)].
\]
左端 \(z=1\) 使 \(A=0\)，右端 \(z_L=1+\alpha L\) 使 \(q\ln z_L=n\pi,\ n=1,2,\ldots\)。于是
\[
\omega_n=\alpha\sqrt{\frac{T_0}{\mu}}
\sqrt{\frac14+\left(\frac{n\pi}{\ln(1+\alpha L)}\right)^2},
\quad X_n(x)\propto(1+\alpha x)^{-1/2}
\sin\frac{n\pi\ln(1+\alpha x)}{\ln(1+\alpha L)}.
\]

(3) 用 \(X\) 乘本征方程，分部积分并使用 \(X(0)=X(L)=0\)，得
\[
\mu\omega^2\int_0^L|X|^2dx=\int_0^L T(x)|X'|^2dx>0.
\]
因此无非零零频或虚频本征模；对两个不同本征值的解交叉相乘并相减，边界项为零，得到 \(\int_0^L\mu X_mX_n\,dx=0\)，即质量加权正交，本题 \(\mu\) 为常量。

(4) 当 \(\alpha L\ll1\)，\(\ln(1+\alpha L)=\alpha L+O(\alpha^2L^2)\)，故
\[
\omega_n=\frac{n\pi}{L}\sqrt{T_0/\mu}[1+O(\alpha L)]
\]
并且 \(X_n(x)\to\sin(n\pi x/L)\)。由于变张力在 \(\alpha\to0\) 时逐点趋向 \(T_0\)，此极限验证了所得谱的正常化。
